\section{Axiomatic frame and readout semantics: minimal structure in HPA--$\Omega$}
\label{sec:axioms_readout}

\subsection{$\Omega$ axioms (O1--O6) and the ``no external time'' starting point}
\label{subsec:axioms}

We adopt the following axioms as a minimal commitment \cite{Ma2025FoP,Ma2025HPOT}:
\begin{itemize}
  \item \textbf{O1 ($\Omega$ axiom).} The theory is specified by a single global state $\omega_{\Omega}$ (not an ensemble); external time is not fundamental.
  \item \textbf{O2 (finite information / holographic bound).} The distinguishable dimension of any finite region is bounded by an exponential of its boundary area (holographic scaling) \cite{tHooft1993,Susskind1995,Bousso2002}.
  \item \textbf{O3 (causal locality).} Local algebras remain finitely propagating under discrete-step evolution; correlation functions are determined by $\omega_{\Omega}$ and a sequence of local automorphisms.
  \item \textbf{O4 (holographic map / QEC structure).} The bulk--boundary map is approximately isometric on a code subspace and supports entanglement-wedge reconstruction (operator-algebraic QEC semantics) \cite{AlmheiriDongHarlow2015,Pastawski2015,Harlow2017}.
  \item \textbf{O5 (scan--projection readout).} Observer ``time'' arises from a scan orbit combined with finite-resolution projection; readout probabilities are given by an effective state $\rho_{\mathrm{eff}}$ and a POVM/effects family $E_k^{(\varepsilon)}$ \cite{Kraus1983,NielsenChuang2000}.
  \item \textbf{O6 (Weyl scan algebra).} Scan shift and phase multiplication form a Weyl pair $(U,V)$ satisfying
  \begin{equation}
    UV = \e^{2\pi \iu \alpha} VU ,
    \qquad \alpha\notin\QQ,
  \end{equation}
  implying intrinsic complementarity (non-simultaneous diagonalizability); see e.g.\ the irrational rotation algebra literature \cite{Rieffel1981}.
\end{itemize}

\noindent These axioms can be read as a minimal operational contract. O1 fixes the ontic ``whole'' (a single state) and removes external time as a primitive. O2 makes resolution a physical constraint rather than an analyst's convenience. O3--O4 provide the structural inputs needed for a local-to-global story: locality in the algebraic evolution and a bulk--boundary encoding consistent with QEC semantics. O5 states that the observer's timeline is not the ontic time but a \emph{scan--projection composite}; O6 fixes the irreducible complementarity of this composite via Weyl noncommutativity.

Operationally, these axioms force any readout to confront a triad:
\begin{equation}
  \text{projection} \;\wedge\; \text{complementarity} \;\wedge\; \text{finite information}.
\end{equation}
CAP will be formulated as a minimal variational closure of the mismatch induced by this triad.

\subsection{Readout operators: irrational rotation, windows, and mechanical words}
\label{subsec:readout_operators}

In the minimal model, the scan is an irrational rotation on the circle:
\begin{equation}
  x_k = x_0 + k\alpha \pmod{1},
  \qquad \alpha\notin\QQ,\quad x_0\in[0,1).
\end{equation}
For irrational $\alpha$, the orbit is equidistributed mod $1$ (Weyl's theorem) \cite{Weyl1916}, providing the canonical ``uniform ontic measure'' that discrepancy quantifies at finite readout depth.
Given a binary window $W\subset[0,1)$ (e.g.\ an interval), the readout sequence is the mechanical word
\begin{equation}
  s_k = \ind_{W}(x_k) \in \{0,1\}.
\end{equation}
For an irrational rotation, binary partitions generate Sturmian flows. The resulting symbolic dynamics has minimal complexity (exactly $n+1$ distinct length-$n$ factors), making it the canonical ``least structured'' non-periodic readout stream \cite{MorseHedlund1940,Lothaire2002}. The golden branch yields the Fibonacci word and equips the tick index with canonical Ostrowski numeration, degenerating to Zeckendorf decomposition in the golden case \cite{Ma2025HPA,Zeckendorf1972}. In CAP terms, canonical coding is not an aesthetic choice but an operational compression: it is how a finite observer assigns stable addresses to scan events under limited capacity.

Readout statistics at finite resolution $\varepsilon$ are specified in operational quantum language by a POVM/effects family $\{E_k^{(\varepsilon)}\}_k$ with $\sum_k E_k^{(\varepsilon)}=\mathds{1}$ and
\begin{equation}
  P_k^{(\varepsilon)} = \Tr\!\left(\rho_{\mathrm{eff}}\, E_k^{(\varepsilon)}\right).
\end{equation}
This keeps ``discrete observed events'' inside standard operational semantics while making explicit the dependence on finite resolution. In particular, the same ontic state $\omega_\Omega$ can induce different effective dynamics at different $\varepsilon$ through the coarse-grained instrument $\{E_k^{(\varepsilon)}\}$, which is precisely the layer separation CAP exploits.

\subsection{Orbit trace and finite part: a canonical regularization convention}
\label{subsec:orbit_trace}

Long-time averages in the scan--readout setting are not arbitrary: CAP requires a fixed convention for taking regulated limits from discrete orbits to continuum functionals. We adopt a canonical \emph{orbit trace} / Abel finite-part convention (Convention R1) \cite{Ma2025HPOT,Ma2025HPA}.

Conceptually, regulated readout sums often differ by additive constants that depend on the chosen regularization scheme. In CAP this ambiguity is not ignorable: mismatch is promoted to a source term, and source terms shift physical potentials unless the constant is fixed consistently. Convention R1 fixes the additive constants in a way compatible with the scan semantics (Abel-type regularization along the orbit), ensuring that ``discrepancy accumulation'' admits a well-defined coarse-grained continuum interpretation that can be coupled to covariant field equations.


